\documentclass[11pt,a4paper]{article}
\usepackage[T1]{fontenc}
\usepackage[utf8]{inputenc}
\usepackage{lmodern}
\usepackage{amsmath,amssymb,amsthm,mathtools}
\usepackage{geometry}
\geometry{margin=2.6cm}
\usepackage{booktabs,array}
\usepackage{enumitem}
\usepackage[hidelinks]{hyperref}
\usepackage{xurl}
\hypersetup{pdftitle={Fourteen initial Theta-data on one curve over Q(i)},pdfauthor={Toshihisa Urahigashi}}
\setlist{itemsep=2pt,topsep=4pt}
\raggedbottom
\emergencystretch=3em

\newcommand{\Z}{\mathbb{Z}}
\newcommand{\Q}{\mathbb{Q}}
\newcommand{\F}{\mathbb{F}}
\newcommand{\Fmod}{F_{\mathrm{mod}}}
\newcommand{\Vbadmod}{\mathbb{V}^{\mathrm{bad}}_{\mathrm{mod}}}
\newcommand{\VV}{\underline{\mathbb{V}}}
\newcommand{\ord}{\operatorname{ord}}
\newcommand{\Nm}{\operatorname{N}}
\newcommand{\GL}{\mathrm{GL}}
\newcommand{\SL}{\mathrm{SL}}
\newcommand{\Gal}{\mathrm{Gal}}
\newcommand{\qq}{\underline{\underline{q}}}
\newcommand{\tS}{\textup{[S]}}
\newcommand{\tP}{\textup{[P]}}

\theoremstyle{plain}
\newtheorem{theorem}{Theorem}[section]
\newtheorem{proposition}[theorem]{Proposition}
\newtheorem{lemma}[theorem]{Lemma}
\theoremstyle{definition}
\newtheorem{remark}[theorem]{Remark}

\title{Fourteen initial $\Theta$-data on one curve over $\mathbb{Q}(i)$}
\author{Toshihisa Urahigashi\\\small ORCID 0009-0004-0460-6242\\\small\texttt{123026.urahigashi.toshihisa@gmail.com}}
\date{Version 1.0.0, 6 November 2026\\\small\href{https://doi.org/10.5281/zenodo.23990717}{doi:10.5281/zenodo.23990717}}

\begin{document}
\maketitle

\begin{abstract}
Let $E\colon y^2=x(x-1)(x-a)$ with $a=(41+26i)^{472}$, a curve over $\Q(i)$. We prove that for each prime $\ell$ in
\[
 \Lambda=\{89,97,101,103,107,109,113,127,131,137,139,149,151,157\}
\]
the curve extends to an initial $\Theta$-datum in the sense of \cite[Def.~3.1]{IUT1}, and that this datum satisfies
the hypotheses of \cite[Thm.~1.10]{IUT4}. The field of moduli is $\Q(i)$. Above the split prime $2357$ there is one
place of multiplicative and one of good reduction, with ramification indices $15\ell$ and $1$. The curve and the
fourteen levels come from \cite[\S13]{Second}. That report showed that the levels pass five numerical conditions, and
said explicitly that they had not been verified to be initial $\Theta$-data; the present report gives the verification.
The proof uses the individual $\ell$ only through five integer facts and one Frobenius witness per level. The norm of
$1-a$ has $1592$ digits, and its factorization is not used. Instead, the valuation of $1-a$ at each split place is
bounded through the cyclotomic factor $\Phi_d(\pi)$, where $d$ is the order of $\pi$ modulo the place, and the
$138{,}731$ primes that this bound does not exclude are tested directly; the inert places and the places above $2357$
are treated directly. The fourteen levels are the primes in
$[\sqrt{\log\mathfrak q^{\forall}},\,472/3)$, and they satisfy condition~(C1) of \cite[Cor.~2.2]{IUT4}. Nothing is
claimed about \cite[Cor.~3.12]{IUT3} or about $abc$.
\end{abstract}

\section{Introduction}\label{s:intro}

\tS{} An initial $\Theta$-datum \cite[Def.~3.1]{IUT1} is a list of objects --- a number field $F$, a once-punctured
elliptic curve $X_F$ over $F$, a prime $\ell$, an orbicurve $\underline C_K$, a set $\VV$ of places of a field $K$, a
set $\Vbadmod$ of places of the field of moduli, and a cusp $\underline\epsilon$ --- subject to conditions (a)--(f).
Dupuy and Hilado work out the curve 11a1 over $\Q$
\cite[\S2.10]{DH1}, and Zhou describes constructions of $\mu_6$-initial $\Theta$-data over $\Q$
\cite[Prop.~2.5, Lemma~2.6]{Zhou1}. Over quadratic fields, \cite{First} gives a datum over $\Q(\sqrt5)$ and
\cite{Second} one over $\Q(i)$, both at $\ell=7$, and the catalogue \cite{Cat} lists these and several others.

\tS{} In \cite[\S13]{Second} the curve of this report appears with the fourteen levels of $\Lambda$, and two parts of
condition~(c) are settled there: $\ell\nmid\Nm(1-a)$, and the sufficient condition for the order condition that
$\Nm(1-a)$ is free of $89$th powers. The same section says that the fourteen levels are \emph{not} verified initial
$\Theta$-data: the argument for the Galois image had been carried out only for the datum at $\ell=7$, and conditions
(a)--(f) had not been checked at these levels. In \cite{Cat} the curve is listed as $D_{2357}^{(\ell)}$, with its values marked as
conditional on such a verification. \tP{} This report supplies it:

\begin{itemize}
\item the image condition at each level, from a transvection at the bad place above $2357$ and one Frobenius witness
 per level (\S\ref{ss:image});
\item the order condition at every selected place, by an argument that does not use the factorization of $1-a$ and
 that is independent of the computation of \cite[\S13]{Second} (Proposition~\ref{p:arith}(d));
\item the remaining conditions of \cite[Def.~3.1]{IUT1} and the hypotheses of \cite[Thm.~1.10]{IUT4}, by the arguments
 of \cite[\S\S3--4]{Second} with $\ell$ in place of $7$ (\S\ref{s:proof}).
\end{itemize}

\tP{} \emph{Scope.} This report constructs data and verifies their defining conditions. Beyond the local invariants
of Theorem~\ref{t:main}(iii) and the margin quoted in Remark~\ref{r:window}, it does not evaluate quantities of
\cite[Cor.~3.12]{IUT3} or \cite[Thm.~1.10]{IUT4} at these data, and it makes no claim about either statement or about
$abc$. Each statement is marked \tS{} if it is a statement of a cited source, paraphrased with its
locator, and \tP{} if it is proved or computed here.

\section{The datum and the theorem}\label{s:datum}

\begin{equation}\label{e:datum}
\begin{aligned}
 &L=\Q(i), &&\pi=41+26i,\quad \Nm(\pi)=2357,\quad a=\pi^{472},\\
 &E\colon y^2=x(x-1)(x-a), &&\ell\in\Lambda,\\
 &F=L(E[15]), &&K=F(E[\ell])=L(E[15\ell]).
\end{aligned}
\end{equation}
Let $X_F=E_F\setminus\{0\}$ and $C_F=[X_F/\{\pm1\}]$. Take as $\Vbadmod$ the set of places of $L$ of odd residue
characteristic at which $E$ has multiplicative reduction. By \S\ref{ss:ab} below, these are the place $v_\pi$ and
the places of odd residue characteristic dividing $1-a$; by Proposition~\ref{p:arith}(e) none of them lies above
$\ell$, so the same set is obtained by taking the places of multiplicative reduction outside $2\ell$, as in \cite[\S2]{Note} and \cite[\S3]{Cat}.

\begin{theorem}\label{t:main}
\tP{} Let $\ell\in\Lambda$.
\begin{enumerate}[label=(\roman*)]
\item The data \eqref{e:datum} extend to an initial $\Theta$-datum
 $(\overline F/F,X_F,\ell,\underline C_K,\VV,\Vbadmod,\underline\epsilon)$ in the sense of \cite[Def.~3.1]{IUT1}, with
 $\Fmod=L$.
\item This datum satisfies the hypotheses of \cite[Thm.~1.10]{IUT4}: $E_F$ has good reduction at every
 nonarchimedean place of $F$ that lies over $\mathbb{V}_{\mathrm{mod}}\setminus\Vbadmod$ and does not divide $2\ell$, and
 $F=\Fmod(\sqrt{-1},E_{\Fmod}[2\cdot3\cdot5])$ in the sense stated there.
\item Above $2357$ there are exactly two places of $L$, the place $v_\pi\in\Vbadmod$ and the place $v_{\bar\pi}$ of good
 reduction. At the places of $\VV$ above them, the ramification indices of $K$ over $\Q_{2357}$ are
 \[
  e_{\mathrm b}=15\ell,\qquad e_{\mathrm g}=1,
 \]
 and $\ord_{2357}(\qq_{\mathrm b})=472/\ell$, where $\qq_{\mathrm b}$ is the $2\ell$-th root of the $q$-parameter of
 \cite[Example~3.2(iv)]{IUT1}. In particular $e\le2357-2=2355$ at both places, and $e_{\mathrm b}=2355$ exactly
 when $\ell=157$.
\end{enumerate}
\end{theorem}

\begin{remark}[one curve, fourteen levels]\label{r:uniform}
\tP{} The field of moduli, the curve, the prime $2357$ and the set $\Vbadmod$ do not depend on $\ell$; only $\ell$ and,
with it, $K$, the covering, $\VV$ and the cusp change. The proof of (i) and (ii) uses the individual level only
through the five integer facts
\[
 \ell\ne2357,\qquad \ell\nmid944,\qquad \ell\nmid\Nm(1-a),\qquad \ell\nmid23040,\qquad \ord_{\mathfrak p}(1-a)\le88<\ell,
\]
and through one Frobenius witness from Table~\ref{tab:witness}; part (iii) also uses $\ell\le157$. The set $\Lambda$ is taken from \cite[\S13]{Second},
where it is the set of primes in a window; the window plays no role in the proof (see \S\ref{s:window}).
\end{remark}

\section{The arithmetic of \texorpdfstring{$a$}{a} and \texorpdfstring{$1-a$}{1-a}}\label{s:arith}

\tP{} The order condition at the bad places dividing $1-a$ rests on the following elementary bound.

\begin{lemma}\label{l:cyclo}
\tP{} Let $\mathfrak q$ be a prime of $\Z[i]$ with $\mathfrak q\ne(\pi)$, whose residue characteristic $p$ is neither
$2$ nor $59$, and suppose that $\mathfrak q\mid\pi^{472}-1$. Let $d$ be the order of $\pi$ in $(\Z[i]/\mathfrak
q)^\times$. Then $d\mid472$, $d\mid\Nm\mathfrak q-1$, and
\[
 \ord_{\mathfrak q}(\pi^{472}-1)=\ord_{\mathfrak q}\bigl(\Phi_d(\pi)\bigr).
\]
In particular $(\Nm\mathfrak q)^k\le\Nm(\Phi_d(\pi))$ for $k=\ord_{\mathfrak q}(\pi^{472}-1)$.
\end{lemma}

\begin{proof}
The group $(\Z[i]/\mathfrak q)^\times$ has order $\Nm\mathfrak q-1$, and $\pi^{472}\equiv1$, so $d$ divides both. Put
$x=\pi^d$ and $m=472/d$. Then $\pi^{472}-1=(x-1)(1+x+\dots+x^{m-1})$, and the second factor is congruent to $m$ modulo
$\mathfrak q$, a unit because $m$ divides $2^3\cdot59$ and $p\nmid2\cdot59$. So $\ord_{\mathfrak
q}(\pi^{472}-1)=\ord_{\mathfrak q}(\pi^d-1)$. Next, $\pi^d-1=\prod_{e\mid d}\Phi_e(\pi)$, and for $e\mid d$, $e<d$, the
factor $\Phi_e(\pi)$ divides $\pi^e-1$, which is prime to $\mathfrak q$ by the minimality of $d$. So
$\ord_{\mathfrak q}(\pi^d-1)=\ord_{\mathfrak q}(\Phi_d(\pi))$. Finally, if $\mathfrak q^k$ divides the nonzero element
$\Phi_d(\pi)$, then $(\Nm\mathfrak q)^k$ divides $\Nm(\Phi_d(\pi))$.
\end{proof}

\begin{table}[ht]
\centering\small
\begin{tabular}{@{}rrrrr@{}}
\toprule
$d$ & $\deg\Phi_d$ & digits of $\Nm(\Phi_d(\pi))$ & $R_d=\lfloor\Nm(\Phi_d(\pi))^{1/89}\rfloor$ & $\mathrm{lcm}(4,d)$\\
\midrule
$1,2,4,8$ & $1,1,2,4$ & $4,4,7,14$ & $1$ & ---\\
$59$, $118$ & $58$ & $196$ & $157$ & $236$\\
$236$ & $116$ & $392$ & $24856$ & $236$\\
$472$ & $232$ & $783$ & $617831325$ & $472$\\
\bottomrule
\end{tabular}
\caption{The cyclotomic factors $\Phi_d(\pi)$, $d\mid472$. Each is computed from $\pi$ by exact division,
$\Phi_d(\pi)=\prod_{e\mid d}(\pi^e-1)^{\mu(d/e)}$, and their product is checked to equal $\pi^{472}-1$.}
\label{tab:cyclo}
\end{table}

\begin{proposition}\label{p:arith}
\tP{}
\begin{enumerate}[label=(\alph*)]
\item $j(E)=256(a^2-a+1)^3/\bigl(a^2(a-1)^2\bigr)$ is not real. Hence $j(E)\notin\Q$ and $\Fmod=\Q(j(E))=L$.
\item $2357$ is prime, $2357=\pi\bar\pi$ splits in $L$, and
 \[
  \ord_\pi(a)=472,\qquad\ord_{\bar\pi}(a)=0,\qquad\ord_\pi(1-a)=\ord_{\bar\pi}(1-a)=0.
 \]
\item $\Nm(1-a)$ is an integer of $1592$ digits with $v_2=9$, $v_3=2$, $v_5=2$, and
 $\gcd\bigl(\mathrm{Re}(1-a),\mathrm{Im}(1-a)\bigr)=240$.
\item At every place $\mathfrak p$ of $L$ of odd residue characteristic, $\ord_{\mathfrak p}(1-a)\le88$.
\item For every $\ell\in\Lambda$: $\ell\ne2357$, and $\ell$ divides none of $944=2^4\cdot59$,
 $23040=2^9\cdot3^2\cdot5$ and $\Nm(1-a)$.
\end{enumerate}
\end{proposition}

\begin{proof}
(a) Write $j=A/B$ with $A=256(a^2-a+1)^3$ and $B=a^2(a-1)^2$ in $\Z[i]$. The imaginary part of $A\bar B$ is a nonzero
integer (of $7953$ digits after removing the common factor with $\Nm(B)$), so $j\ne\bar j$. As a second check, at the
two places of $L$ above $29$ ($i\equiv12$ and $i\equiv17$) the residue of $a$ is $20$ and $23$, and $j$ reduces to $16$
and $18$ in $\F_{29}$; a rational $j$ would reduce to the same element at both. The field of moduli of an elliptic
curve is generated by its $j$-invariant \cite[Prop.~1.8(ii)]{IUT4}, and $j\in L\setminus\Q$.

(b) $\bar\pi$ is not an associate of $\pi$. Modulo $\pi$ one has $a\equiv0$, so $1-a\equiv1$. Modulo $\bar\pi$ one has
$i\equiv41\cdot26^{-1}$ and $\pi\equiv82$; since $82^4\equiv202\not\equiv1\pmod{2357}$ and $\gcd(472,2356)=4$,
$82^{472}\not\equiv1$, so $\bar\pi\nmid1-a$.

(c), (e) Exact computation (\S\ref{s:repro}); $\Nm(1-a)\bmod\ell$ was also recomputed in $\F_\ell[i]$ from $\pi$.

(d) If $\mathfrak p=(p)$ is inert, then $\ord_{\mathfrak p}(1-a)$ is the smaller of $v_p(\mathrm{Re}(1-a))$ and
$v_p(\mathrm{Im}(1-a))$, so $\ord_{\mathfrak p}(1-a)\le v_p(240)\le1$; this covers $p=3$ and $p=59$. The places above $2357$ are covered by (b). The remaining places are
the places $\mathfrak q$ of degree one above the primes $p\equiv1\pmod4$, $p\ne2357$, and for them we use
Lemma~\ref{l:cyclo} and Table~\ref{tab:cyclo}. Suppose $\ord_{\mathfrak q}(1-a)\ge89$, and let $d$ be the order of
$\pi$ modulo $\mathfrak q$. By the lemma, $p=\Nm\mathfrak q\le R_d=\lfloor\Nm(\Phi_d(\pi))^{1/89}\rfloor$, and
$p\equiv1\pmod{\mathrm{lcm}(4,d)}$. For $d\mid8$ this is impossible, since $R_d=1$. For $d\in\{59,118\}$ it is
impossible as well, since $R_d=157$ and no $p\equiv1\pmod{236}$ is at most $157$. So $p$ is a prime $p\equiv1\pmod{236}$
with $p\le24856$ or a prime $p\equiv1\pmod{472}$ with $p\le617{,}831{,}325$. There are $138{,}731$ such primes, and at
each of their $277{,}462$ places the valuation of $1-a$ was computed. Exactly four of these places divide $1-a$, one
above each of $p=1181$, $21713$, $2876369$ and $6188393$, and each divides it exactly once. Hence no place has
$\ord_{\mathfrak q}(1-a)\ge89$.
\end{proof}

\begin{remark}[a second computation]\label{r:second}
\tS{} \cite[\S13]{Second} proves the stronger statement that $\Nm(1-a)$ is free of $89$th powers. It splits
$\Phi_{472}$ over $\Z[i]$ into two factors of degree $116$, sieves each of the nine resulting block norms up to its own
$89$th root (at most $24856$; $8318$ primes in all), and checks with the pairwise greatest common divisors that only
$2$ and $5$ divide two blocks. Its program is in the deposit of \cite{Second}. \tP{} The computation above does not split
$\Phi_{472}$ and shares no code with it; both run in seconds and agree. A direct search for $89$th powers in
$\Nm(1-a)$ would need the primes up to $\Nm(1-a)^{1/89}\approx10^{17.9}$.
\end{remark}

\section{Proof of Theorem~\ref{t:main}}\label{s:proof}

\subsection{Conditions (a) and (b)}\label{ss:ab}

\tS{} Condition~(a) of \cite[Def.~3.1]{IUT1} asks that $F$ be a number field with $\sqrt{-1}\in F$. Condition~(b)
asks that $X_F$ have stable reduction at every nonarchimedean place of $F$; that $\Vbadmod$ be a nonempty set of
nonarchimedean places of $\Fmod$ of odd residue characteristic such that $X_F$ has bad multiplicative reduction at
every place of $F$ above $\Vbadmod$; that $F/\Fmod$ be Galois of degree prime to $\ell$; and that the $2\cdot3$-torsion
points of $E_F$ be rational over $F$.

\tP{} $i\in L\subseteq F$. The field of moduli of $X_F$ is that of $E$, which is $L$ by
Proposition~\ref{p:arith}(a). The extension $F=L(E[15])$ is Galois over $L$, and its group embeds in
$\GL_2(\Z/15\Z)$, of order $48\cdot480=23040=2^9\cdot3^2\cdot5$, so $[F:L]$ is prime to $\ell$
(Proposition~\ref{p:arith}(e)). The $2$-torsion of $E$ is rational over $L$ (Legendre form), and the $3$-torsion is
rational over $F$ since $E[3]\subseteq E[15]$.

\tS{} By \cite[Prop.~1.8(v)]{IUT4}, if $k$ is a perfect complete discretely valued field, $l\ge3$ is a prime invertible in its
ring of integers and all $l$-torsion points of an elliptic curve $E_k$ are $k$-rational, then $E_k$ has semistable
reduction. \tP{} At a completion $F_w$ apply this with $l=3$ if the residue characteristic is not $3$, and with $l=5$
otherwise. So $E_F$ has semistable reduction at every nonarchimedean place: good where $\ord(j)\ge0$ and
multiplicative where $\ord(j)<0$ \cite[Props.~VII.5.1, VII.5.5]{Sil}. At a place of good reduction the minimal model is
smooth, and at a place of multiplicative reduction its special fibre is a nodal cubic on which the origin is a smooth
point; in both cases $X_F$ has stable reduction.

\tP{} The discriminant of the model \eqref{e:datum} is $16a^2(a-1)^2$, and $c_4=16(a^2-a+1)$. At an odd place where $a$
and $1-a$ are units, the discriminant is a unit and $E$ has good reduction. At $v_\pi$, and at an odd place $\mathfrak p$
dividing $1-a$, one has $a^2-a+1\equiv1$, so $c_4$ is a unit, the model is minimal \cite[\S VII.1]{Sil}, and the reduction is
multiplicative,
with
\[
 \ord_\pi(j)=-944,\qquad \ord_{\mathfrak p}(j)=-2\ord_{\mathfrak p}(1-a).
\]
So $\Vbadmod$ consists of $v_\pi$ and the places of odd residue characteristic dividing $1-a$; it is nonempty, and its
elements have odd residue characteristic. Since $\ord(j)<0$ and the reduction over $F$ is semistable, the reduction is
multiplicative at every place of $F$ above $\Vbadmod$.

\subsection{Condition (c): residue characteristics and orders}\label{ss:orders}

\tS{} Condition~(c) asks that $\ell\ge5$ be a prime; that the image of the outer homomorphism $G_F\to\GL_2(\F_\ell)$
determined by the $\ell$-torsion points of $E_F$ contain $\SL_2(\F_\ell)$, $K$ being the field cut out by its kernel;
and that $\ell$ be prime to the residue characteristics of the places in $\Vbadmod$ and to the orders of the
$q$-parameters of $E_F$ at the places of $F$ above $\Vbadmod$ (the ``local heights'' of \cite[Def.~3.3]{GenEll}).

\tP{} The residue characteristics of $\Vbadmod$ are $2357$ and the odd primes dividing $\Nm(1-a)$, and $\ell$ is none of
them (Proposition~\ref{p:arith}(e)). The orders of the $q$-parameter are $\ord_\pi(q)=-\ord_\pi(j)=944$ at $v_\pi$ and
$\ord_{\mathfrak p}(q)=2\ord_{\mathfrak p}(1-a)$ at an odd place $\mathfrak p\mid1-a$, where
$1\le\ord_{\mathfrak p}(1-a)\le88$ by Proposition~\ref{p:arith}(d). At a place $w$ of $F$ above $v$ the order is
multiplied by $e(w/v)$, which divides $[F:L]\mid23040$. As $\ell$ is a prime with $\ell\nmid944$, $\ell\nmid23040$ and
$\ell\ge89$, it divides none of these orders.

\subsection{Condition (c): the image}\label{ss:image}

\begin{lemma}\label{l:transvection}
\tP{} Let $\ell$ be a prime and let $H\subseteq G\subseteq\GL_2(\F_\ell)$ be subgroups with $H$ normal in $G$. Suppose
that $H$ contains a transvection $T\ne1$ and that $G$ contains an element $g$ whose characteristic polynomial is
irreducible over $\F_\ell$. Then $H\supseteq\SL_2(\F_\ell)$.
\end{lemma}

\begin{proof}
Let $D=\ker(T-1)$, a line. Since the characteristic polynomial of $g$ has no root in $\F_\ell$, $g$ has no eigenvector,
so $gD\ne D$. Then $T'=gTg^{-1}$, which lies in $H$, is a transvection $\ne1$ with fixed line $gD$. In a basis
$e_1\in D$, $e_2\in gD$ the two read $\bigl(\begin{smallmatrix}1&x\\0&1\end{smallmatrix}\bigr)$ and
$\bigl(\begin{smallmatrix}1&0\\y&1\end{smallmatrix}\bigr)$ with $x,y\ne0$. Since $\ell$ is prime, their powers are all
the matrices $\bigl(\begin{smallmatrix}1&t\\0&1\end{smallmatrix}\bigr)$ and
$\bigl(\begin{smallmatrix}1&0\\t&1\end{smallmatrix}\bigr)$, $t\in\F_\ell$, and these generate $\SL_2(\F_\ell)$.
\end{proof}

\tP{} We apply the lemma with $H$ and $G$ the images of $\Gal(\overline\Q/F)$ and $\Gal(\overline\Q/L)$ in
$\GL(E[\ell])\simeq\GL_2(\F_\ell)$; $H$ is normal in $G$ because $F/L$ is Galois.

\tP{} \emph{The transvection.} Let $w$ be a place of $F$ above $v_\pi$. Over $F_w$ the curve is a split Tate curve
(\S\ref{ss:split}; \cite[App.~C, Thm.~14.1]{Sil}), with $\ord_w(q)=e(w/v_\pi)\cdot944$, and this is prime to $\ell$ by
\S\ref{ss:orders}. As $\ell\ne2357$, the maximal unramified extension of $F_w$ contains $\mu_\ell$, and $q^{1/\ell}$
generates a ramified extension of it of degree $\ell$. So an element of the inertia group at $w$ acts on
$E[\ell]=\langle\zeta_\ell,q^{1/\ell}\rangle$ by a transvection $T\ne1$, and $T\in H$.

\tP{} \emph{The Frobenius witnesses.} For a prime $q\equiv1\pmod4$ and a root $r$ of $r^2\equiv-1\pmod q$, the place
$\mathfrak q=(q,i-r)$ of $L$ sends $a$ to $\mathrm{Re}(a)+r\,\mathrm{Im}(a)\bmod q$. At the four places of
Table~\ref{tab:witness} this residue is neither $0$ nor $1$, so $E$ has good reduction there. The point count gives
$a_{\mathfrak q}=q+1-\#E(\F_q)$. For $\ell\ne q$, a Frobenius element $g\in G$ at $\mathfrak q$ acts on
$E[\ell]\simeq\tilde E[\ell]$ \cite[Prop.~VII.3.1]{Sil} with characteristic polynomial $X^2-a_{\mathfrak q}X+q$
\cite[Thm.~V.2.3.1]{Sil}, which is irreducible over $\F_\ell$ exactly when $a_{\mathfrak q}^2-4q$ is a nonsquare
modulo $\ell$.

\begin{table}[ht]
\centering\small
\begin{tabular}{@{}rrrrrrl@{}}
\toprule
$q$ & $i\equiv$ & $a\equiv$ & $\#E(\F_q)$ & $a_{\mathfrak q}$ & $a_{\mathfrak q}^2-4q$ & levels covered\\
\midrule
13 & 5  & 3  & 16 & $-2$ & $-48=-3\cdot4^2$  & 89, 101, 107, 113, 131, 137, 149\\
17 & 13 & 16 & 16 & $2$  & $-64=-8^2$        & 103, 107, 127, 131, 139, 151\\
29 & 12 & 20 & 32 & $-2$ & $-112=-7\cdot4^2$ & 89, 97, 101, 103, 131, 139, 157\\
97 & 75 & 75 & 96 & $2$  & $-384=-6\cdot8^2$ & 89, 109, 113, 137, 139, 157\\
\bottomrule
\end{tabular}
\caption{Frobenius witnesses. The last column lists the $\ell\in\Lambda$, $\ell\ne q$, at which $a_{\mathfrak q}^2-4q$
is a nonsquare modulo $\ell$. Every $\ell\in\Lambda$ occurs in at least one row.}
\label{tab:witness}
\end{table}

\tP{} By quadratic reciprocity, the first three rows are the conditions $\ell\equiv2\pmod3$, $\ell\equiv3\pmod4$ and
$(\ell/7)=-1$. The fourth, $(-6/\ell)=-1$, is needed only at $\ell=109$, which no other row covers. The choice of place
matters: above $17$ the other place ($i\equiv4$) sends $a$ to $1$ and belongs to $\Vbadmod$, and above $97$ the other
place ($i\equiv22$) has $\#E(\F_{97})=112$, $a_{\mathfrak q}=-14$ and discriminant $-192$.

\tP{} So, for every $\ell\in\Lambda$, Lemma~\ref{l:transvection} applies with the transvection at $w$ and a witness
from Table~\ref{tab:witness}, and the image of $\Gal(\overline\Q/F)$ contains $\SL_2(\F_\ell)$.

\subsection{Split reduction, and the local invariants above 2357}\label{ss:split}

\tP{} At an odd place where $a\equiv0$, the special fibre of \eqref{e:datum} is $y^2=x^2(x-1)$, with tangent cone
$y^2=-x^2$ at the node; it splits because $i\in L$. At an odd place where $a\equiv1$, putting $t=x-1$ gives
$y^2=(1+t)t^2$, with tangent cone $y^2=t^2$, already split. So the reduction is split multiplicative over $L$ at every
place of $\Vbadmod$, and it stays split over $F$ and $K$.

\tP{} At $v_{\bar\pi}$ the reduction is good and $15\ell$ is prime to $2357$, so $K$ is unramified there and
$e_{\mathrm g}=1$. At $v_\pi$ the curve is a split Tate curve over $\Q_{2357}$ with $\ord_\pi(q)=944$. Adjoining
$E[15\ell]$ adjoins $\mu_{15\ell}$, which is unramified, and a $15\ell$-th root of $q$; since $\gcd(15\ell,944)=1$, the
ramification index is exactly $15\ell$. So $e_{\mathrm b}=15\ell$ and $\ord_{2357}(\qq_{\mathrm b})=944/(2\ell)=472/\ell$.
Finally $15\ell\le15\cdot157=2355=2357-2$, with equality only at $\ell=157$.

\subsection{Condition (d): the quotient and the core}\label{ss:core}

\tP{} Write $V=E[\ell](K)$, fix a line $W\subset V$ and an element $0\ne\xi\in Q=V/W$, and let
$\pi_W\colon E_K\to E'=E_K/W$, with $\psi\colon E'\to E_K$ such that $\psi\pi_W=[\ell]$. Restriction gives an
isomorphism $\bar\pi_W\colon Q\simeq\ker\psi$, so this kernel is a constant group over $K$. Put $\underline
X_K=E'\setminus\ker\psi$, $\underline C_K=[\underline X_K/\{\pm1\}]$ and $C_K=[X_K/\{\pm1\}]$. The covering $\underline
X_K\to X_K$ is Galois with group $\ker\psi\simeq Q$. It extends to the isogeny $\psi$ of the compactifications, so it is
unramified above the cusp, and its fibre there, $\ker\psi$, consists of $K$-rational points; so the decomposition group
of the cusp maps trivially to $Q$. Passing to the quotients by $\pm1$, with the cusp $0$ of $\underline X_K$ as origin, gives an orbicurve of type
$(1,\ell\text{-tors})_\pm$ in the sense of \cite[Def.~2.1]{EtTh}. Let $\underline\epsilon$ be the cusp corresponding to
$\pm\bar\pi_W(\xi)$.

\tS{} By \cite[Prop.~2.1]{MFHMP}, if the once-punctured curve of an elliptic curve over a field of characteristic zero is
arithmetic, then its $j$-invariant is one of $2^{14}\cdot31^3\cdot5^{-3}$, $2^2\cdot73^3\cdot3^{-4}$, $1728$, $0$.
\tP{} Here $j(E)$ is not real (Proposition~\ref{p:arith}(a)), so $X_{\overline K}$ is not arithmetic. The finite \'etale
map $X_{\overline K}\to C_{\overline K}$ puts the two orbicurves in the same commensurability class, so $C_{\overline K}$ is
not arithmetic either. \tS{} A punctured hemi-elliptic orbicurve over an algebraically closed field of characteristic
zero that is not arithmetic is a core \cite[Prop.~2.7]{CanLift}; an orbicurve over $k$ is a $k$-core if and only if its
base change to a separable closure is a core there \cite[Prop.~2.3(i)]{CanLift}; and the property is unchanged under
extension of algebraically closed fields of characteristic zero \cite[Prop.~2.3(ii)]{CanLift}. \tP{} So $C_K$ is a
$K$-core, and, since $\underline C_K\to C_K$ is finite \'etale, it is the $K$-core of $\underline C_K$ required by
\cite[Def.~3.1(d)]{IUT1}. For a completion $K_v$, extend $K\hookrightarrow\overline K$ to an embedding of $\overline K$
into an algebraic closure of $K_v$ and apply \cite[Prop.~2.3(ii)]{CanLift} and then \cite[Prop.~2.3(i)]{CanLift}; no
real completion occurs, since $i\in K$. The open subgroups $\Delta_{\underline{\underline X}}\subseteq\Delta_{\underline{\underline
C}}\subseteq\Delta_{\underline C}$ of \cite[Def.~3.1(d)]{IUT1} are supplied by \cite[Prop.~2.2, Def.~2.3]{EtTh}. Since
$\ord_\pi(j)=-944<0$, $j$ is not an algebraic integer and $E$ has no complex multiplication; this is not used above.

\subsection{Conditions (e) and (f): the places and the cusp}\label{ss:ef}

\tS{} Condition~(e) asks for a subset $\VV\subseteq\mathbb V(K)$ mapping bijectively to $\mathbb V(\Fmod)$, such that
at every $\underline v\in\VV$ above $\Vbadmod$ the orbicurve $\underline C_{\underline v}$ is of type
$(1,\Z/\ell\Z)_\pm$ \cite[Def.~2.5(i)]{EtTh}. That definition makes two assumptions: that the quotient
$\overline\Pi{}^{\mathrm{ell}}_X\twoheadrightarrow Q$ factor through the natural quotient $\Pi_X\twoheadrightarrow\widehat\Z$
determined by the quotient $\Pi_X^{\mathrm{tp}}\twoheadrightarrow\Z$ of \cite[\S1]{EtTh}, and that the splitting of $\overline D_x\twoheadrightarrow G_K$ which determines the covering
$\underline{\underline X}\to\underline X$ be compatible with the $\{\pm1\}$-structure of \cite[Thm.~1.10(iii)]{EtTh},
which is defined when the residue characteristic is odd. Condition~(e) records that the $2$-torsion part of (b) gives the condition
``$K=\ddot K$'' of \cite[Def.~2.5(i)]{EtTh}, and that the models of type $(1,(\Z/\ell\Z)^\Theta)$ and
$(1,(\Z/\ell\Z)^\Theta)_\pm$ at such places are supplied by the theory of \cite[\S2]{EtTh}. Condition~(f) asks that
$\underline\epsilon$ arise from a nonzero element of $Q$, and that at each $\underline v\in\VV$ above $\Vbadmod$ its
localization arise, up to sign, from the canonical generator of the quotient $\widehat\Z$ of \cite[Def.~2.5(i)]{EtTh}.

\tP{} At each $u\in\Vbadmod$ choose a place of $F$ above $u$ and a place $w$ of $K$ above it. The split Tate sequence
$0\to\mu_\ell\to E(K_w)[\ell]\to\Z/\ell\to0$, in which the class of $q^{1/\ell}$ maps to $1$, determines a line
$W_w=\mu_\ell\subset V$ and an element $\pm\xi_w$ of $V/W_w$, up to sign: changing the root $q^{1/\ell}$ changes it by an
element of $\mu_\ell=W_w$, which does not change its class, and reversing the orientation of the dual graph changes the
sign. The group $\SL(V)$ acts transitively on such pairs: for a pair $(W',\pm\xi')$ choose $s\in W'$ and a lift $t$ of
$\xi'$ with $\det(s,t)=1$, and note that $\SL(V)$ acts transitively on the bases so obtained. So, by \S\ref{ss:image},
there is
$\sigma_u\in\Gal(K/F)$ with $\sigma_u(W_w,\pm\xi_w)=(W,\pm\xi)$. Put $v(u)=\sigma_uw$, the place with
$|b|_{\sigma_uw}=|\sigma_u^{-1}b|_w$; the isomorphism $\sigma_u^{-1}\colon K_{v(u)}\simeq K_w$ commutes with
localization, so the fixed global pair $(W,\pm\xi)$ acquires at $v(u)$ the local marking. The global $W$ and
$\underline\epsilon$ are not changed from place to place. Since $\sigma_u$ fixes $F$, the place $v(u)$ lies above the
same $u$, and the fibres above different $u$ are disjoint, so these choices can be made simultaneously. Choosing an
arbitrary place of $K$ above each remaining place of $\Fmod$ gives $\VV\xrightarrow{\ \sim\ }\mathbb V(\Fmod)$.

\tP{} Locally at $v(u)$ the quotient is $E(q)\to E(q^\ell)$, $z\mapsto z^\ell$, and its dual is $E(q^\ell)\to E(q)$,
$z\mapsto z$, which multiplies the period of the dual graph by $\ell$. So the quotient $\Pi_X\to Q$ factors through
$\widehat\Z\to\Z/\ell$, which is the first assumption of \cite[Def.~2.5(i)]{EtTh}. The second concerns only the
splitting that determines $\underline{\underline X}_{v(u)}$; we take it to be the one of the models supplied by
\cite[\S2]{EtTh} and recorded in condition~(e). So $\underline C_{v(u)}$ is of type $(1,\Z/\ell\Z)_\pm$. Since $\ker\psi=\{q^r\bmod
q^{\ell\Z}:r\in\Z/\ell\}$ and $\pi_W(q^{1/\ell})=q$, the localization of $\underline\epsilon$ is the cusp of the
generator $\pm1$: the label counts deck translations by the period $q$, and no Tate order or ramification index enters
it. The rational $2$-torsion gives ``$K=\ddot K$'', and the Weil pairings give $\mu_{60}\subset F$ and $\mu_\ell\subset
K$.

\subsection{The hypotheses of \texorpdfstring{\cite[Thm.~1.10]{IUT4}}{IUT IV, Thm. 1.10}}\label{ss:thm110}

\tS{} \cite[Thm.~1.10]{IUT4} is stated for initial $\Theta$-data in the situation of \cite[Cor.~3.12]{IUT3}. It
assumes that $E_F$ has good reduction at every nonarchimedean place of $F$ over
$\mathbb V_{\mathrm{mod}}\setminus\Vbadmod$ that does not divide $2\ell$, and that $F=\Fmod(\sqrt{-1},E_{\Fmod}[2\cdot3\cdot5])$;
that is, $F$ is obtained from the field $F_{\mathrm{tpd}}=\Fmod(E_{\Fmod}[2])$ by adjoining $\sqrt{-1}$ and the
$(3\cdot5)$-torsion of the model of $E$ over $F_{\mathrm{tpd}}$ given by a Legendre equation.

\tP{} Every place of $L$ of odd residue characteristic at which $E$ has bad reduction is in $\Vbadmod$
(\S\ref{ss:ab}), so the first assumption holds. For the second, $E[2]\subseteq E(L)$, so
$F_{\mathrm{tpd}}=L$; $\sqrt{-1}\in L$; \eqref{e:datum} is a Legendre equation over $L$; and $F=L(E[15])$.

\section{The levels and condition (C1) of \texorpdfstring{\cite[Cor.~2.2]{IUT4}}{IUT IV, Cor. 2.2}}\label{s:window}

\tS{} \cite[Cor.~2.2(i)]{IUT4} defines $\log(\mathfrak q^\forall)$ as the normalized degree of the effective arithmetic
divisor determined by the $q$-parameters of the curve at all nonarchimedean places. The initial $\Theta$-data
constructed in \cite[Cor.~2.2(ii)]{IUT4} have a prime $\ell$ with
\[
 \textup{(C1)}\qquad (\log\mathfrak q^\forall)^{1/2}\le\ell\le10\delta\cdot(\log\mathfrak q^\forall)^{1/2}\cdot
 \log(2\delta\cdot\log\mathfrak q^\forall),\qquad \delta=2^{12}\cdot3^3\cdot5\cdot d,
\]
where $d$ bounds the degree of the point $x_E\in U_X(\overline\Q)^{\le d}$ of the $\lambda$-line.

\tP{} Over $F$ the curve is semistable, and at a place $w$ of multiplicative reduction $\ord_w(q)=-\ord_w(j)$. Since
the normalized degree does not change under extension of the base field \cite[Def.~1.9(i)]{IUT4}, the sum can be taken
over the nonarchimedean places $v$ of $L$:
$\log\mathfrak q^\forall=\tfrac12\sum_v\max(0,-\ord_vj)\log\Nm v$. The contributions are $944$ at $v_\pi$,
$2\ord_{\mathfrak p}(1-a)$ at an odd place $\mathfrak p\mid1-a$, and $2$ at the place above $2$, where
$\ord_{1+i}(j)=3\cdot8-(8+2\cdot9)=-2$. So
\begin{align*}
 \log\mathfrak q^\forall&=472\log2357+\log\Nm(1-a)-8\log2=7324.751610927\ldots,\\
 (\log\mathfrak q^\forall)^{1/2}&=85.584762726\ldots,
\end{align*}
and $\Lambda$ is exactly the set of primes in $[(\log\mathfrak q^\forall)^{1/2},\,472/3)$. With $d=2=[\Q(a):\Q]$ the
upper bound of (C1) is about $2.2251\times10^{10}$, and it only grows with $d$. So every $\ell\in\Lambda$ satisfies
(C1).

\begin{remark}\label{r:window}
\tP{} The upper end $472/3$ of the window is not part of (C1) or of \cite[Def.~3.1]{IUT1}. \tS{} It is the threshold
$\ell<N/3$ of \cite[\S\S8,13]{Second}, where $N=472$, at which the margin $\ord_{2357}(\qq_{\mathrm b})-3=(472-3\ell)/\ell$
of the first stage is still positive; it is $1/157$ at $\ell=157$. \cite[Cor.~2.2]{IUT4} also requires that $x_E$ lie in
a compactly bounded subset satisfying the condition $(\ast_{j\text{-inv}})$ at $2$ and outside a finite exceptional set
$\mathfrak{Exc}_d$, which is not given explicitly. \tP{} These requirements are not checked here; (C1) is the only condition
of \cite[Cor.~2.2]{IUT4} compared with the data.
\end{remark}

\section{Reproduction}\label{s:repro}

\begin{itemize}
\item In the archive deposited with this note, \texttt{scripts/d2357\_check.py} (Python~3, standard library only,
 about ten seconds) recomputes the numbers of Proposition~\ref{p:arith} --- for part~(d), the valuations at the
 $277{,}462$ places of the $138{,}731$ primes; the reduction to these places is Lemma~\ref{l:cyclo} and the case
 analysis in the proof --- as well as Table~\ref{tab:cyclo}, Table~\ref{tab:witness} and the remarks after it, the local
 invariants of \S\ref{ss:split} and the values of \S\ref{s:window}, and compares them with the numbers printed here. It
 also compares the reduction of $a$ modulo every odd prime below $3000$, $p\ne2357$, with direct modular
 exponentiation. Its output is \texttt{outputs/d2357\_check.txt}; it exits with a nonzero status if a comparison
 fails.
\item The second computation of Remark~\ref{r:second} is \texttt{scripts/powerfree\_certificate.py} in the deposit of
 \cite{Second}; run unchanged, it reproduces the counts stated there.
\item The arguments that are not computations --- Lemma~\ref{l:cyclo} and the reduction in the proof of
 Proposition~\ref{p:arith}(d), the image argument of \S\ref{ss:image} apart from the point counts and the residues, the
 stable and the split reduction, the core, the choice of places and cusp, and \S\ref{ss:thm110} --- are not checked by
 these programs.
\end{itemize}

\section*{Use of generative AI}
Generative AI tools (large language models) were used extensively in this work, under the author's direction, in
finding the example and its proofs, in the computations and in the writing. The author takes full responsibility for
the contents of this report.

\begin{thebibliography}{MFHMP}\raggedright
\bibitem[CanLift]{CanLift}
Shinichi Mochizuki, \emph{The absolute anabelian geometry of canonical curves}, Doc. Math., Extra Volume Kato (2003),
609--640.

\bibitem[Cat]{Cat}
Toshihisa Urahigashi, \emph{Explicit initial $\Theta$-data and the two readings of Step~(v) of {IUT}~{IV},
{T}heorem~1.10: a catalogue}, version 1.0.0, Zenodo, 2 November 2026,
\href{https://doi.org/10.5281/zenodo.23870448}{doi:10.5281/zenodo.23870448}.

\bibitem[DH1]{DH1}
Taylor Dupuy and Anton Hilado, \emph{The statement of {M}ochizuki's {C}orollary 3.12: initial theta data}, Ramanujan J.
\textbf{68} (2025), no.~1, Paper No.~18, \href{https://doi.org/10.1007/s11139-025-01100-y}{doi:10.1007/s11139-025-01100-y};
arXiv:2004.13228v2 [math.NT]. The section number cited refers to the arXiv version.

\bibitem[EtTh]{EtTh}
Shinichi Mochizuki, \emph{The \'etale theta function and its {F}robenioid-theoretic manifestations}, Publ. Res. Inst.
Math. Sci. \textbf{45} (2009), no.~1, 227--349, \href{https://doi.org/10.2977/prims/1234361159}{doi:10.2977/prims/1234361159}.

\bibitem[First]{First}
Toshihisa Urahigashi, \emph{An explicit initial $\Theta$-datum with mixed reduction at a split prime}, version 0.1.7,
Zenodo, \href{https://doi.org/10.5281/zenodo.22537342}{doi:10.5281/zenodo.22537342}.

\bibitem[GenEll]{GenEll}
Shinichi Mochizuki, \emph{Arithmetic elliptic curves in general position}, Math. J. Okayama Univ. \textbf{52} (2010),
1--28.

\bibitem[IUT1]{IUT1}
Shinichi Mochizuki, \emph{Inter-universal {T}eichm\"uller theory {I}: construction of {H}odge theaters}, Publ. Res.
Inst. Math. Sci. \textbf{57} (2021), 3--207, \href{https://doi.org/10.4171/PRIMS/57-1-1}{doi:10.4171/PRIMS/57-1-1}.

\bibitem[IUT3]{IUT3}
Shinichi Mochizuki, \emph{Inter-universal {T}eichm\"uller theory {III}: canonical splittings of the log-theta-lattice},
Publ. Res. Inst. Math. Sci. \textbf{57} (2021), 403--626, \href{https://doi.org/10.4171/PRIMS/57-1-3}{doi:10.4171/PRIMS/57-1-3}.

\bibitem[IUT4]{IUT4}
Shinichi Mochizuki, \emph{Inter-universal {T}eichm\"uller theory {IV}: log-volume computations and set-theoretic
foundations}, Publ. Res. Inst. Math. Sci. \textbf{57} (2021), 627--723, \href{https://doi.org/10.4171/PRIMS/57-1-4}{doi:10.4171/PRIMS/57-1-4}.

\bibitem[MFHMP]{MFHMP}
Shinichi Mochizuki, Ivan Fesenko, Yuichiro Hoshi, Arata Minamide and Wojciech Porowski, \emph{Explicit estimates in
inter-universal {T}eichm\"uller theory}, Kodai Math. J. \textbf{45} (2022), no.~2, 175--236, \href{https://doi.org/10.2996/kmj45201}{doi:10.2996/kmj45201}.

\bibitem[Note]{Note}
Toshihisa Urahigashi, \emph{A local non-containment in Step~(v) of {IUT}~{IV}, {T}heorem~1.10}, version 3.0.0, Zenodo,
1 November 2026, \href{https://doi.org/10.5281/zenodo.23852752}{doi:10.5281/zenodo.23852752}.

\bibitem[Second]{Second}
Toshihisa Urahigashi, \emph{An initial $\Theta$-datum over an imaginary quadratic field, and numerical values of the
local quantities of inter-universal {T}eichm\"uller theory~{IV}}, version 1.0.5, Zenodo,
\href{https://doi.org/10.5281/zenodo.22601333}{doi:10.5281/zenodo.22601333}.

\bibitem[Sil]{Sil}
Joseph H. Silverman, \emph{The arithmetic of elliptic curves}, second edition, Graduate Texts in Mathematics
\textbf{106}, Springer, Dordrecht, 2009, \href{https://doi.org/10.1007/978-0-387-09494-6}{doi:10.1007/978-0-387-09494-6}.

\bibitem[Zhou1]{Zhou1}
Zhong-Peng Zhou, \emph{The inter-universal {T}eichm\"uller theory and new {D}iophantine results over the rational
numbers.~{I}}, arXiv:2503.14510 [math.NT], 2025.
\end{thebibliography}

\end{document}
